\section{The endowment theorem}\label{sec:endo}

This section derives the EP model's master equation, exhibits the canonical
feasible point, maps the feasibility region over the full parameter grid,
and proves the compression identity that ties the model to the independent
TMR4A estimate. Throughout, ``feasible'' means
$|\pi - \pi_{\mathrm{obs}}| \le 15\%$ with $\pi_{\mathrm{obs}} = 10^{-3}$.

\subsection{Transmission of a four-allele endowment}\label{sec:trans}
The first question is mechanical: can the couple's four allelic states at a
locus survive reproduction at all? Let the father carry alleles $\{A, B\}$
and the mother $\{C, D\}$; each child independently receives one paternal
and one maternal allele.

\begin{proposition}[Four-allele transmission]\label{prop:trans}
The probability that all four founder alleles are represented among $n$
children is
\begin{equation}\label{eq:p4}
P_4(n) \;=\; \bigl[1 - 2^{1-n}\bigr]^2 ,
\end{equation}
so that $P_4(2) = 0.250$, $P_4(3) = 0.563$, $P_4(4) = 0.766$,
$P_4(5) = 0.879$, and $P_4(10) = 0.996$.
\end{proposition}

\begin{proof}
The events ``both paternal alleles appear'' and ``both maternal alleles
appear'' are independent. A specific paternal allele is absent from all $n$
children with probability $2^{-n}$, and the two absence events are
mutually exclusive, so the probability that both paternal alleles appear is
$1 - 2^{1-n}$; likewise for the maternal pair. Multiplying gives
Eq.~\eqref{eq:p4}.
\end{proof}

The values $0.25$, $0.56$, and $0.77$ for $n = 2, 3, 4$ follow directly
from Eq.~\eqref{eq:p4}, and the rapid saturation is the quantitative form
of the founder-flush intuition that rapid population growth means that few
variants in the original couple are lost. Two consequences matter for what
follows. First,
the expansion phase of the EP model is not a near-loss event for the
endowment: with $\gamma \ge 1.6$ the family size reaches double digits
within a few generations and per-locus transmission failures become rare.
Second, the losses that do occur are precisely priced by the hazard
integral \eqref{eq:hazard}, which is why the master equation below needs no
separate retention factor beyond $e^{-\Lambda_c}$.

\subsection{The master equation of the endowed-pair model}\label{sec:master}
Consider a pair of gene copies sampled today and traced backward. It
coalesces during the post-bottleneck history (plateau and expansion) with
the standard hazard \eqref{eq:hazard}; the density of coalescence at
backward time $t$ is $\lambda(t)e^{-\Lambda(t)}$ with
$\lambda(t) = 1/(2N(t))$, contributing the expected depth
\begin{equation}\label{eq:Dbelow}
D_{<c} \;=\; \int_0^{T_0-1} t\,\lambda(t)\,e^{-\Lambda(t)}\,\mathrm{d}t .
\end{equation}
A pair that survives to the couple generation (probability
$e^{-\Lambda_c}$, with $\Lambda_c$ the hazard accumulated to the couple's
generation) passes through the four founder copies. With probability
$\tfrac14$ the pair lands on the same founder copy and coalesces at $T_0$;
with probability $\tfrac34$ it lands on two different copies. For the
different-copy landings, the two copies are unrelated \emph{within the
model} --- their divergence is the endowment $h$ itself, independent of
whether the two copies belong to the same founder (divergence $=$ that
founder's own heterozygosity $\approx h$) or to different founders
(divergence $d_{AE} \approx (h_A + h_E)/2 \approx h$ for disjoint,
comparable founders). The mixture therefore collapses to a single
endowment depth $h/(2\mu)$ added below the couple:

\begin{proposition}[EP master equation]\label{prop:master}
Under an endowed-pair trajectory $E(T_0, \gamma, \Ne^{\max}, h)$,
\begin{equation}\label{eq:master}
\boxed{\;
E[T_2] \;=\; D_{<c} \;+\; e^{-\Lambda_c}\Big[T_0 + \tfrac{3}{4}\,\frac{h}{2\mu}\Big],
\qquad \pi \;=\; 2\mu\,E[T_2],
\;}
\end{equation}
where $D_{<c}$ and $\Lambda_c$ are computed from the demographic function
by Eqs.~\eqref{eq:Dbelow} and \eqref{eq:hazard}.
\end{proposition}

\begin{proof}
Immediate from the decomposition above: a pair either coalesces before the
couple (contributing $D_{<c}$ in expectation through the hazard-weighted
integral \eqref{eq:Dbelow}) or reaches the couple generation, where the
same-copy case contributes $T_0$ and the different-copy case contributes
$T_0 + h/(2\mu)$. Weighting the latter two cases by $\tfrac14$ and
$\tfrac34$ gives the bracket in Eq.~\eqref{eq:master}. The heterozygosity
form of the same statement --- present-day diversity
$\pi = 2\mu\,D_{<c} + H_{\mathrm{pool}}\,e^{-\Lambda_c}$ with
$H_{\mathrm{pool}} = \tfrac34 h$ --- follows by multiplying the bracket by
$2\mu$ and agrees with the drift-decay view in which the pool heterozygosity
is multiplied by the exact retention product
$\prod_t(1 - 1/2N(t)) = e^{-\Lambda_c}$.
\end{proof}

Equation~\eqref{eq:master} is the EP analogue of the sampled-couple
equation \eqref{eq:sampled}, with the term $\tfrac34\,(2\Nepre)$ replaced by
$\tfrac34\,h/(2\mu)$. The comparison is the paper's central observation in
equation form: the depth that the sampled-couple model imports from a
pre-bottleneck population of size $\Nepre$ is imported by the EP model from
the couple's own divergence. For the canonical founders,
\begin{equation}\label{eq:depth}
\frac{h}{2\mu} \;=\; \frac{1.576\times10^{-3}}{2.4\times10^{-8}}
\;=\; 6.57\times10^{4}\ \text{generations}
\;=\; 1.64\ \text{Myr},
\end{equation}
and $\tfrac34$ of this --- $4.92\times10^{4}$ generations $= 1.23$ Myr ---
is the endowment's weighted contribution to $E[T_2]$: more than the
\emph{entire} required $E[T_2] = 4.17\times10^{4}$ of constraint
\eqref{eq:T2req}. The endowment
is not a marginal patch; it is a depth reservoir of the same magnitude as
the constraint.

\subsection{The generation bound}\label{sec:bound}
The endowment is a transmission requirement, and the master equation prices
it exactly. Inverting Eq.~\eqref{eq:master} at $\pi = \pi_{\mathrm{obs}}$,
\begin{equation}\label{eq:hreq}
h \;=\; \frac{4}{3}\,\frac{\pi_{\mathrm{obs}} \;-\;
2\mu\bigl[D_{<c} + e^{-\Lambda_c}T_0\bigr]}{e^{-\Lambda_c}} ,
\end{equation}
and since pairs coalescing before the couple have coalescence times below
$T_0$, so that $D_{<c} \le (1-e^{-\Lambda_c})\,T_0$, while
$e^{-\Lambda_c} \le 1$, every trajectory of the EP class obeys
\begin{equation}\label{eq:floor}
\pi \;\le\; 2\mu T_0 \;+\; \tfrac{3}{4}\,h
\qquad\Longrightarrow\qquad
h \;\ge\; \tfrac{4}{3}\,\bigl(\pi_{\mathrm{obs}} - 2\mu T_0\bigr)
\;=\; 0.63\times10^{-3},
\end{equation}
or $1.9$ million SNVs per founder. The interpretation is the theorem's
affirmative content, stated as a bound. Mutation accumulated below the
couple can contribute at most $2\mu T_0 = 0.53\times10^{-3}$ per site ---
$53\%$ of $\pi_{\mathrm{obs}}$, or $1.6$ million SNVs per founder --- no
matter how fast the population grows; the remaining $47\%$ at minimum must
be \emph{carried} into the founding event on the couple's four copies. At
the canonical trajectory the split is $44.7\%$ de novo and $55.3\%$ carried
(Eq.~\eqref{eq:canonval}): the endowment backbone
$0.75\,h\,e^{-\Lambda_c} = 0.553\times10^{-3}$ exceeds the entire de-novo
component. The historical objection --- that two individuals cannot supply
the observed variation --- is correct in exactly this sense and no
stronger: two individuals cannot \emph{generate} population-scale diversity
in $550$ kyr, and the theorem never claims otherwise; what two primates
\emph{can} do is carry it, and Eq.~\eqref{eq:floor} prices the carrying.
It is worth naming the identity that connects the two model classes: the
canonical endowment's depth $h/(2\mu) = 6.57\times10^{4}$ generations
equals $2\Ne$ at $\Ne = 3.3\times10^{4}$ --- numerically the
$\Nepre = 3\times10^{4}$ of the sampled-couple class. The replacement of
$\tfrac34(2\Nepre)$ by $\tfrac34\,h/(2\mu)$ is therefore exact not only in
form but in magnitude; what the EP class changes is not the depth of the
reservoir but its \emph{assignment} --- from an ancestry to an initial
condition --- and Section~\ref{sec:nature} returns to what follows from
that and what does not.

\subsection{The canonical point}\label{sec:canonical}
The canonical parameterization of the paper is
\begin{equation}\label{eq:canonical}
T_0 = 22{,}000\ \text{generations}\ (550\ \text{kya}), \quad
\gamma = 1.6, \quad
\Ne^{\max} = 3\times10^{4}, \quad
h = 1.576\times10^{-3},
\end{equation}
the endowment corresponding to $4.73$ million SNVs per founder --- $116\%$
of the 1000 Genomes average and $83\%$ of the chimpanzee value
\citep{PradoMartinez2013,Auton2015}. Evaluated exactly (fixed-seed script,
Section~\ref{sec:sims}): $\Lambda_c = 0.759$,
$e^{-\Lambda_c} = 0.468$, $D_{<c} = 8{,}311$ generations, giving
\begin{equation}\label{eq:canonval}
E[T_2] \;=\; 8{,}311 + 0.468\times\bigl[22{,}000 + 49{,}250\bigr]
\;=\; 41{,}671\ \text{generations},
\qquad
\pi \;=\; 1.000\times10^{-3}.
\end{equation}
The master constraint \eqref{eq:T2req} is met to within $0.01\%$; the
expected pairwise coalescence depth, $1{,}042$ kyr, lies inside the
observed autosomal range of roughly $0.5$--$2.0$ Myr
\citep{GarriganHammer2005,LiDurbin2011}. Three independent checks
corroborate the decomposition at the same point. The coalescent simulation
of the de-novo component (mutations younger than $T_0$, msprime,
Section~\ref{sec:sims}) yields $\pi_{\mathrm{denovo}} = 0.448\times10^{-3}$,
matching the master equation's $2\mu\,[D_{<c} + e^{-\Lambda_c}T_0] =
0.447\times10^{-3}$ to $0.3\%$. The endowment backbone, evaluated exactly
as $H_{\mathrm{pool}}\,e^{-\Lambda_c} = 0.75 \times 1.576\times10^{-3}
\times 0.468 = 0.553\times10^{-3}$, brings the total to
$1.000\times10^{-3}$ --- the same figure the master equation produces,
now assembled from two independent simulation routes.

\subsection{The feasibility region}\label{sec:feasregion}
Figure~\ref{fig:reqendo} displays the inversion of Eq.~\eqref{eq:master}:
for each $\Ne^{\max}$, the endowment $h_{\mathrm{req}}(T_0)$ required to
reach $\pi_{\mathrm{obs}}$. The horizontal lines mark the biological
bounds \eqref{eq:hbounds}. The figure carries the feasibility claim
graphically: at $\Ne^{\max} = 3\times10^{4}$, the required endowment lies
\emph{between} the human-average and chimpanzee bounds across the entire
range $T_0 \in [400, 600]$ kya; at $\Ne^{\max} = 10^{4}$ it exceeds the
chimpanzee bound at every $T_0$, so the canonical long-term
$\Ne \approx 10^{4}$ is \emph{not} sufficient for sole progenitorship ---
a genuine restriction of the EP class, reported here as such and discussed
in Section~\ref{sec:limits}.

\begin{figure}[t]
\centering
\includegraphics[width=0.72\textwidth]{../outputs/fig1_required_endowment.pdf}
\caption{Required endowment $h_{\mathrm{req}}(T_0)$ for
$\pi_{\mathrm{obs}} = 10^{-3}$ at $\gamma = 1.6$, for three post-expansion
plateau sizes. Horizontal lines: human-average endowment ($4.1$M SNVs,
dotted) and chimpanzee endowment ($5.7$M SNVs, dashed)
\citep{PradoMartinez2013,Auton2015}. Shaded band: the joint feasibility
window of Section~\ref{sec:window}. At $\Ne^{\max} = 3\times10^4$ the
required endowment lies within the primate range across the window; at
$\Ne^{\max} = 10^4$ it does not.}
\label{fig:reqendo}
\end{figure}

Table~\ref{tab:gridextract} extracts the $540$-point feasibility grid
($12$ values of $T_0 \times 5$ plateau sizes $\times 3$ growth rates
$\times 3$ endowments) at the canonical growth rate and endowment. The
plateau-size requirement is the grid's sharpest message:
$\Ne^{\max} \ge 2\times10^{4}$ is necessary, and values of
$2$--$5\times10^{4}$ --- consistent with the mid-Pleistocene estimates at
the upper range of PSMC reconstructions and with the transient
$\sim10^{5}$ peaks inferred by MSMC \citep{LiDurbin2011,SchiffelsDurbin2014}
--- give feasibility across the whole window.

\begin{table}[t]
\centering\small
\caption{Feasible $T_0$ ranges (kya) at $\gamma = 1.6$ and the canonical
endowment $h = 1.576\times10^{-3}$, by plateau size
($|\pi - 10^{-3}| \le 15\%$). Full grid: \texttt{feasibility\_grid.csv}.}
\label{tab:gridextract}
\begin{tabular}{lccccc}
\toprule
$\Ne^{\max}$ & $10^4$ & $2\times10^4$ & $3\times10^4$ & $5\times10^4$ & $10^5$ \\
\midrule
feasible $T_0$ (kya) & none & 300--800 & 150--800 & 150--600 & 150--450 \\
\bottomrule
\end{tabular}
\end{table}

\begin{figure}[t]
\centering
\includegraphics[width=0.72\textwidth]{../outputs/fig3_feasibility_map.pdf}
\caption{Nucleotide diversity $\pi$ (in units of $10^{-3}$) over the
$(T_0, \Ne^{\max})$ plane at $\gamma = 1.6$ and the canonical endowment.
The white-outlined band is the joint feasibility window
(Section~\ref{sec:window}). The $10^{-3}$ level band runs through it at
$\Ne^{\max} \ge 2\times10^4$.}
\label{fig:feasmap}
\end{figure}

\subsection{TMR4A and the compression identity}\label{sec:window}
Under sole progenitorship every locus in every living individual passes
through the couple's four gene copies: at time $T_0$ the global ancestry
contains exactly four lineages per locus, never more. The time at which the
genome-wide ancestry compresses to four lineages --- the quantity
\citet{Swamidass2019} named TMR4A and estimated from reconstructed
ancestral recombination graphs --- is therefore \emph{identically} the
couple's time:

\begin{proposition}[Compression identity]\label{prop:tmr4a}
Under an EP trajectory, $\mathrm{TMR4A}_{\mathrm{true}} = T_0$ exactly, for
every parameterization.
\end{proposition}

\begin{proof}
At $T_0$ the population consists of two diploid individuals, i.e.\ four
gene copies per locus. No locus can have more than four ancestral lineages
at $T_0$, and (absent a pre-bottleneck population) none has fewer to which
the four extant clades could compress earlier. The genome-wide compression
time --- the infimum of times at which every locus has $\le 4$ lineages ---
is therefore $T_0$.
\end{proof}

The direct-clock estimate of this quantity is $495 \pm 100$ kya
\citep{Swamidass2019,Swamidass2022}, obtained from autosomal coalescent
reconstructions with pedigree-calibrated rates and recombination maps
--- that is, on the \emph{direct} clock that Section~\ref{sec:framework}
singled out as calibration-free. Setting
$\mathrm{TMR4A}_{\mathrm{true}} = T_0$ restricts
$T_0 \in [395, 595]$ kya. Intersecting with the diversity-feasible region
of Table~\ref{tab:gridextract} gives the \textbf{joint feasibility window}
\begin{equation}\label{eq:window}
T_0^{*} \in [400,\,595]\ \text{kya},
\qquad \text{nominal } T_0^{*} = 550\ \text{kya},
\end{equation}
represented on the paper's grid by the feasible points $\{400, 450, 500,
550\}$ kya at $\gamma = 1.6$, $\Ne^{\max} = 2$--$3\times10^4$,
$h = 1.576\times10^{-3}$. It is worth pausing on what has happened here.
The independent autosomal argument of Swamidass (four-lineage compression)
and the backward-simulation feasibility analysis of \citet{HossjerGauger2019}
(a single-couple origin as recent as $\sim$500 kya consistent with genomic
data) already converged on the $400$--$600$ kya window for a \emph{bottleneck
of two}; the diversity constraint that in the sampled-couple class pulls the
bottleneck date toward the present is absent from the EP class, because the
depth it demanded is supplied by the endowment. The arguments now meet in
the same window from three independent directions --- compression,
backward simulation, and pairwise diversity. And the modern spectrum-based
literature has reached the same qualitative verdict: SFS-based inference
recognizes a severe bottleneck in human ancestry \citep{Hu2023}, confirming
that the observed spectrum is bottleneck-shaped and relocating the debate
from \emph{whether} to \emph{when} --- exactly the question the EP window
answers.

Two consistency checks close the section. First, the model's predicted
per-locus coalescence-depth distribution tops out at
$T_0 + h/(2\mu) \approx 550 + 1{,}642 \approx 2.2$ Myr for cross-founder
pairs, consistent with the deepest observed human--human coalescences
(~$2$--$3$ Myr at the deepest loci
\citep{GarriganHammer2005,LiDurbin2011}) once locus-to-locus variance in
the endowment is included. Second, the expected pairwise depth at the
canonical point, $1{,}042$ kyr, is the quantity PSMC-type methods read as
the ``typical'' coalescence depth of the genome; the observed value
($\pi_{\mathrm{obs}}/2\mu \approx 1.04$ Myr) is matched by construction
--- this is Eq.~\eqref{eq:T2req} satisfied, not assumed.
